%-------------------------------------------------------------------------------
%	SECTION TITLE
%-------------------------------------------------------------------------------
\cvsection{论文}

\begin{cvparagraph}
  我是一名自学成才的研究者，受到Yin Wang和Todd Becker工作的启发，经过三年的实验，撰写了一篇关于逆转近视和自然视力恢复的学术风格论文，并发表了一篇关于从零训练GPT-2处理SEC文件的AI研究论文。在计算机科学领域，我仍在努力取得类似的突破。\end{cvparagraph}

%-------------------------------------------------------------------------------%
%	CONTENT
%-------------------------------------------------------------------------------%
\begin{cventries}

%---------------------------------------------------------%
  \cventry
    {李智维} % Author
    {自然视力恢复：实验验证、散光减少与"刚好清晰"原则} % Paper title
    {李智维的博客} % Published in
    {2026年8月} % Date
    {
      \begin{cvitems}
        \item {将三篇论文合并为一项统一研究：验证了自然视力恢复方法，分析了逆转近视过程中的散光减少，并引入了"刚好清晰"原则。}
        \item {在线发布: \href{https://lzwjava.com/natural-vision-restoration-en}{https://lzwjava.com/natural-vision-restoration-en}}
      \end{cvitems}
    }

%---------------------------------------------------------%
  \cventry
    {李智维} % Author
    {SEC-EDGAR-GPT: 基于SEC EDGAR文件从零训练的GPT-2（124M）语言模型} % Paper title
    {GitHub} % Published in
    {2026年6月} % Date
    {
      \begin{cvitems}
        \item {使用nanoGPT在单个RTX 4070上从零训练1.24亿参数的GPT-2，基于SEC EDGAR财务文件（10-K、10-Q、8-K）的15.5亿token。约8小时收敛到验证损失2.28。}
        \item {在线发布: \href{https://github.com/lzwjava/sec-edgar-gpt/raw/main/sec-edgar-gpt.pdf}{sec-edgar-gpt.pdf}}
        \item {代码: \href{https://github.com/lzwjava/sec-edgar-gpt}{github.com/lzwjava/sec-edgar-gpt} \ | \ 模型: \href{https://huggingface.co/lzwjava/sec-edgar-gpt-124m-hf}{huggingface.co/lzwjava/sec-edgar-gpt-124m-hf}}
      \end{cvitems}
    }

%---------------------------------------------------------
\end{cventries}
